\documentclass[11pt]{article}
\usepackage[utf8]{inputenc}
\usepackage{amsmath,amssymb}
\usepackage{hyperref}
\title{Efficient Cancer Immunotherapy Mechanism: A Resource-Aware Approach}
\author{Anonymous}
\date{}
\begin{document}
\maketitle

\begin{abstract}
This paper studies Cancer Immunotherapy Mechanism. Grounded in a real, citable literature \cite{ref1,ref2,ref3,ref4,ref5,ref6,ref7,ref8},
we propose a focused method and report \textbf{illustrative} experiments.
All references below are real and verifiable (OpenAlex/arXiv); the reported
numbers are simulated to illustrate intended behaviour.
\end{abstract}

\section{Introduction}
Cancer Immunotherapy Mechanism is an active research area. This work targets a concrete gap identified from
the recent literature and contributes a method together with an evaluation protocol.

\section{Related Work (real literature)}
The following works, retrieved from public bibliographic APIs, frame our study:
\begin{itemize}
\item Pardoll (2012) — "The blockade of immune checkpoints in cancer immunotherapy", Nature reviews. Cancer (cited 13837).
\item Ribas et al. (2018) — "Cancer immunotherapy using checkpoint blockade", Science (cited 6820).
\item Jiang et al. (2018) — "Signatures of T cell dysfunction and exclusion predict cancer immunotherapy response", Nature Medicine (cited 5640).
\item Sharma et al. (2017) — "Primary, Adaptive, and Acquired Resistance to Cancer Immunotherapy", Cell (cited 5554).
\item Charoentong et al. (2017) — "Pan-cancer Immunogenomic Analyses Reveal Genotype-Immunophenotype Relationships and Predictors of Response to Checkpoint Blockade", Cell Reports (cited 5431).
\item Dong et al. (2002) — "Tumor-associated B7-H1 promotes T-cell apoptosis: A potential mechanism of immune evasion", Nature Medicine (cited 5044).
\end{itemize}

\section{Method}
We reduce the dominant computational cost via adaptive resource allocation, activating only the components needed per input. We instantiate this idea for Cancer Immunotherapy Mechanism and analyse its cost/quality trade-off.

\section{Experiments (Illustrative)}
\begin{center}
\begin{tabular}{lcc}
System & Primary Metric & Efficiency \\ \hline
Strong baseline & 0.812 & 1.00$\times$ \\
Ablation & 0.828 & 0.92$\times$ \\
\textbf{Ours} & \textbf{0.851} & \textbf{0.71$\times$}
\end{tabular}
\end{center}
\emph{Metrics simulated to illustrate the intended effect; not measured.}

\section{Conclusion}
We connected a real literature to a concrete method for Cancer Immunotherapy Mechanism. Future work will
validate the illustrative results with real experiments.

\begin{thebibliography}{9}
\bibitem{ref1} Drew M. Pardoll. The blockade of immune checkpoints in cancer immunotherapy. \emph{Nature reviews. Cancer}, 2012. DOI:10.1038/nrc3239
\bibitem{ref2} Antoni Ribas, Jedd D. Wolchok. Cancer immunotherapy using checkpoint blockade. \emph{Science}, 2018. DOI:10.1126/science.aar4060
\bibitem{ref3} Peng Jiang, Shengqing Gu, Deng Pan et al.. Signatures of T cell dysfunction and exclusion predict cancer immunotherapy response. \emph{Nature Medicine}, 2018. DOI:10.1038/s41591-018-0136-1
\bibitem{ref4} Padmanee Sharma, Siwen Hu‐Lieskovan, Jennifer A. Wargo et al.. Primary, Adaptive, and Acquired Resistance to Cancer Immunotherapy. \emph{Cell}, 2017. DOI:10.1016/j.cell.2017.01.017
\bibitem{ref5} Pornpimol Charoentong, Francesca Finotello, Mihaela Angelova et al.. Pan-cancer Immunogenomic Analyses Reveal Genotype-Immunophenotype Relationships and Predictors of Response to Checkpoint Blockade. \emph{Cell Reports}, 2017. DOI:10.1016/j.celrep.2016.12.019
\bibitem{ref6} Haidong Dong, Scott E. Strome, Diva R. Salomão et al.. Tumor-associated B7-H1 promotes T-cell apoptosis: A potential mechanism of immune evasion. \emph{Nature Medicine}, 2002. DOI:10.1038/nm730
\bibitem{ref7} Roy S. Herbst, Daniel Morgensztern, Chris Boshoff. The biology and management of non-small cell lung cancer. \emph{Nature}, 2018. DOI:10.1038/nature25183
\bibitem{ref8} Roy Noy, Jeffrey W. Pollard. Tumor-Associated Macrophages: From Mechanisms to Therapy. \emph{Immunity}, 2014. DOI:10.1016/j.immuni.2014.06.010
\end{thebibliography}
\end{document}
